%%%%%%%%%%%%%%%%%%%%%%%%%%%%%%%%%%%%%%%%%%%%%%%%%%%%%%%%%%%%%%%%%%%%%%%%%%%%%%%
% SECTION 11.9: Complexity and Uncertainty Effects
% Axioms: AWX-A46 to AWX-A55
% Appendix UNMAPPED_AWA Reference: Section 3.7
%%%%%%%%%%%%%%%%%%%%%%%%%%%%%%%%%%%%%%%%%%%%%%%%%%%%%%%%%%%%%%%%%%%%%%%%%%%%%%%

The interaction between complexity and uncertainty creates multiplicative 
barriers to awareness. While each factor alone reduces awareness, their 
combination produces effects greater than the sum of parts.

\subsubsection{Complexity Effects (AWX-A46 to A47a)}

Complexity refers to the number of elements, relationships, and contingencies 
that must be mentally processed to understand consequences.

\paragraph{AWX-A46: Complexity Decay.}
Awareness decays as a function of perceived complexity:
\begin{equation}
A_{complex} = A_{base} \cdot e^{-\kappa \cdot K}
\label{eq:complexity-decay}
\end{equation}
where $K \in [0,1]$ measures perceived complexity and $\kappa$ is the 
decay rate (empirically $\kappa \approx 1.5$).

\paragraph{AWX-A47: Framing Effects.}
Identical objective complexity can produce different $A$ levels depending 
on how information is framed:
\begin{equation}
A_{framed} = A_{base} \cdot f(Frame)
\end{equation}

\paragraph{AWX-A47a: Failure Awareness Gap.}
Awareness of potential failures is systematically lower than awareness of 
potential successes, even when objectively equiprobable:
\begin{equation}
A_{failure} < A_{success} \quad \text{when } P(failure) = P(success)
\end{equation}

\subsubsection{Four Types of Uncertainty (AWX-A48 to A51)}

Uncertainty affects awareness through distinct channels:

\begin{table}[htbp]
\centering
\caption{Uncertainty Types and Their Effects on Awareness}
\label{tab:uncertainty-types}
\begin{tabular}{@{}llp{6cm}l@{}}
\toprule
\textbf{Axiom} & \textbf{Type} & \textbf{Source} & \textbf{Effect} \\
\midrule
AWX-A48 & Environmental & External volatility, market fluctuations & $A_{long-term} \downarrow$ \\
AWX-A49 & Interpersonal & Others' behavior unpredictable & $A^{comm} \downarrow$ \\
AWX-A50 & Institutional & Rules unclear, enforcement variable & $A_{SYS2} \downarrow$ \\
AWX-A51 & Epistemic & Knowledge itself uncertain & Belief filter $\uparrow$ \\
\bottomrule
\end{tabular}
\end{table}

\paragraph{AWX-A48: Environmental Uncertainty.}
When external conditions are volatile, awareness of long-term consequences 
diminishes because outcomes feel unpredictable:
\begin{equation}
A_{long-term} = A_{base} \cdot (1 - U_{env})
\end{equation}

\paragraph{AWX-A49: Interpersonal Uncertainty.}
When others' responses cannot be anticipated, community-level awareness 
($A^{comm}$) is reduced:
\begin{equation}
A^{comm} = A^{comm}_{base} \cdot (1 - U_{inter})
\end{equation}

\paragraph{AWX-A50: Institutional Uncertainty.}
Unclear rules reduce the effectiveness of deliberative (SYS 2) processing:
\begin{equation}
A_{SYS2} = A_{SYS2,base} \cdot (1 - U_{inst})
\end{equation}

\paragraph{AWX-A51: Epistemic Uncertainty.}
When knowledge itself is uncertain, the belief filter strengthens, 
making motivated reasoning more likely:
\begin{equation}
\lambda_{belief} = \lambda_{base} \cdot (1 + U_{epist})
\end{equation}

\subsubsection{Uncertainty Dynamics (AWX-A52 to A53)}

\paragraph{AWX-A52: Accumulation.}
Multiple uncertainty sources accumulate, but not linearly:
\begin{equation}
U_{total} = 1 - \prod_{i}(1 - U_i)
\end{equation}

\paragraph{AWX-A53: Differentiation.}
Different uncertainty types require different interventions; reducing 
one type does not automatically reduce others.

\subsubsection{The K×U Interaction (AWX-A54)}

\begin{axiom}[AWX-A54: Complexity × Uncertainty Interaction]
When complexity ($K$) and uncertainty ($U$) co-occur, their effects multiply:
\begin{equation}
A_{adjusted} = A_{base} \cdot (1 - K \times U)
\label{eq:kxu-interaction}
\end{equation}
\end{axiom}

This multiplicative interaction explains why new technologies, policy reforms, 
and organizational changes face disproportionate resistance: they combine 
high complexity with high uncertainty.

\begin{example}[Technology Adoption (CASE-07)]
A new software system with:
\begin{itemize}
    \item Complexity $K = 0.6$ (moderately complex interface)
    \item Uncertainty $U = 0.7$ (unclear whether it will work)
\end{itemize}
Produces:
\[
A_{adjusted} = A_{base} \cdot (1 - 0.6 \times 0.7) = A_{base} \cdot 0.58
\]
A \textbf{42\% penalty} on awareness of the new technology's benefits, 
independent of any rational cost-benefit calculation.
\end{example}

\subsubsection{Institutional Mitigation (AWX-A55)}

\paragraph{AWX-A55: Institutional Help.}
Strong institutions can reduce both complexity and uncertainty, thereby 
increasing awareness:
\begin{equation}
A_{institutional} = A_{base} \cdot (1 + \Psi_I \cdot \gamma_I)
\end{equation}

This creates a positive feedback loop: high institutional trust ($\Psi_I$) 
reduces perceived complexity and uncertainty, which increases awareness, 
which enables better decisions, which further strengthens institutional trust.


%------------------------------------------------------------------------------
\subsubsection{Uncertainty Dampening: AWX-A8a}
%------------------------------------------------------------------------------

Beyond the static K×U interaction, AWX-A8a captures a \textbf{dynamic effect}:
uncertainty slows the \emph{rate} at which awareness builds over time.

\begin{axiom}[AWX-A8a: Unsicherheits-Dämpfung]
Risk and ambiguity slow awareness buildup; higher uncertainty leads to lower 
learning rates.

\textbf{Formal Specification:}
\begin{equation}
\frac{dA}{dt} \downarrow \text{ as } U \uparrow
\end{equation}
\end{axiom}

\paragraph{Mechanism.}
Under uncertainty, individuals:
\begin{itemize}
    \item Discount new information more heavily
    \item Require more evidence before updating beliefs
    \item Show slower convergence to stable awareness levels
    \item May exhibit oscillating awareness patterns (approach-avoidance)
\end{itemize}

\paragraph{Example: Climate Change.}
The scientific uncertainty around precise climate outcomes (even if direction 
is clear) dampens the rate at which public awareness builds. Each new study 
updates awareness only marginally, because the base uncertainty remains high.
This explains the multi-decade timescale of awareness development despite 
decades of scientific communication.

\subsubsection{Implications}

The K×U interaction has profound implications for behavioral interventions:

\begin{enumerate}
    \item \textbf{Information alone fails}: Providing more information 
          increases complexity, which can \emph{reduce} awareness
    \item \textbf{Simplification helps}: Reducing perceived complexity 
          directly increases awareness
    \item \textbf{Certainty signals matter}: Reducing uncertainty 
          (through guarantees, trials, social proof) increases awareness
    \item \textbf{Sequence matters}: Address uncertainty before complexity, 
          as uncertainty amplifies complexity effects
\end{enumerate}
